% !TEX root = ../Tesis.tex
\chapter*{Resumen}

La contaminación atmosférica constituye uno de los principales desafíos 
ambientales asociados al crecimiento urbano e industrial. La expansión 
de las actividades industriales y el incremento en la demanda energética 
han favorecido el uso de recursos energéticos fósiles, cuya combustión representa 
una fuente importante de emisiones contaminantes y gases de efecto invernadero. 
Asimismo, este proceso genera vapor de agua ($\mathrm{H_2O_vap}$), que
por su volumen de emisión puede modificar la composición isotópica 
atmosférica local y puede proporcionar información sobre la 
influencia o intensidad de las actividades antropogénicas en la atmósfera.

El vapor de agua derivado de la combustión de recursos energéticos fósiles presenta 
una firma isotópica característica, 
empobrecidos en $\delta\mathrm{D}$, enriquecidos en $\delta^{18}\mathrm{O}$ y con
valores altamente negativos de \textit{d-excess} (=$\delta\mathrm{D}$ - $\delta^{18}\mathrm{O}$). 
En este contexto, la presente 
investigación evalúa la contribución del vapor de agua generado por la 
combustión de recursos energéticos fósiles a la composición isotópica del vapor 
atmosférico local, como un indicador complementario de la influencia 
antropogénica en la atmósfera.

Se recolectaron muestras de vapor atmosférico en la ciudad de Salamanca, Guanajuato. Es
una de las ciudades con mayor actividad industrial de México e índice de contaminación 
atmosférica (Gobierno del Estado de Guanajuato y Secretaría de Seguridad Pública,
2023). Se realizaron 
dos campañas de muestreo en los meses de febrero-marzo y mayo del 2026. Para estimar la contribución 
del vapor de combustión se implementó un modelo binario de mezcla, 
considerando como miembros extremos la firma isotópica del vapor derivado 
de la combustión de recursos energéticos ($\delta^{18}\mathrm{O}_c = 9$\,\textperthousand{} 
y $\delta\mathrm{D}_c = -134.4$\,\textperthousand{}; Xing et al.,2020) y la línea de agua 
meteórica regional ($\delta\mathrm{D} = 7.349\,\delta^{18}\mathrm{O} + 0.693$; Ramírez-González et al., 2024). 
Adicionalmente, se realizaron simulaciones lagrangianas de dispersión 
atmosférica mediante FLEXPART para evaluar la sensibilidad de la relación 
fuente-receptor del CO, como trazador asociado a procesos de combustión.

Los resultados del modelo de mezcla estimaron una contribución de vapor de 
combustión de recursos energéticos entre 18.9 y 21.5\,\%. Asimismo, se identificaron valores 
de \textit{d-excess} de hasta $-38.85$\,\textperthousand{}. Las simulaciones de 
dispersión y el análisis de contaminantes criterio sugieren una influencia 
predominantemente local, con posibles aportaciones de fuentes industriales 
ubicadas al noreste de la ciudad, dentro de su zona metropolitana.

Los resultados evidencian una contribución detectable de vapor de agua 
asociado a la combustión de recursos energéticos fósiles en la composición 
isotópica del vapor atmosférico de Salamanca. Estos hallazgos 
proporcionan elementos para comprender la intensidad de las emisiones 
antropogénicas sobre la variación  
del vapor de agua a escala local. Asimismo, destacan el potencial de los 
isótopos estables del vapor atmosférico como una herramienta complementaria 
para el estudio de la contaminación atmosférica y sus fuentes.

\chapter*{Abstract}

Air pollution is one of the main environmental challenges associated 
with urban and industrial growth. The expansion of industrial activities 
and the increasing energy demand have favored the use of fossil energy 
resources, whose combustion is a major source of pollutant emissions and 
greenhouse gases. This process also generates water vapor ($\mathrm{H_2O_{vap}}$), 
which, owing to its emission volume, can alter the local atmospheric 
isotopic composition and may provide information on the influence 
or intensity of anthropogenic activities in the atmosphere.

Water vapor derived from the combustion of fossil energy 
resources has a characteristic isotopic signature, being depleted 
in $\delta\mathrm{D}$, enriched in $\delta^{18}\mathrm{O}$, 
and showing highly negative \textit{d-excess} values 
(=$\delta\mathrm{D}$ - $\delta^{18}\mathrm{O}$). In this context, 
this study evaluates the contribution of water vapor generated 
by the combustion of fossil energy resources to the isotopic 
composition of local atmospheric vapor, as a complementary 
indicator of anthropogenic influence on the atmosphere.

Atmospheric vapor samples were collected in the city of 
Salamanca, Guanajuato, one of the cities with the highest 
industrial activity and air pollution index in Mexico 
(Gobierno del Estado de Guanajuato y Secretaría de Seguridad Pública, 2023). 
Two sampling campaigns were conducted in February--March and May 2026. To 
estimate the contribution of combustion vapor, a binary mixing model was 
implemented, considering as end-members the isotopic signature of vapor derived 
from the combustion of energy resources ($\delta^{18}\mathrm{O}_c = 9$\,\textperthousand{} 
and $\delta\mathrm{D}_c = -134.4$\,\textperthousand{}; Xing et al., 2020) and 
the regional meteoric water line ($\delta\mathrm{D} = 7.349\,\delta^{18}\mathrm{O} + 0.693$; 
Ramírez-González et al., 2024). Additionally, Lagrangian atmospheric dispersion 
simulations were performed using FLEXPART to evaluate the sensitivity of the 
source--receptor relationship of CO, as a tracer associated with combustion processes.

The mixing model results estimated a combustion vapor contribution from 
energy resources of between 18.9 and 21.5\,\%. In addition, \textit{d-excess} 
values as low as $-38.85$\,\textperthousand{} were identified. The dispersion 
simulations and the analysis of criteria pollutants suggest a predominantly 
local influence, with possible contributions from industrial sources 
located northeast of the city, within its metropolitan area.

The results show a detectable contribution of water vapor associated with 
the combustion of fossil energy resources to the isotopic composition of 
atmospheric vapor in Salamanca. These findings provide insights into the 
intensity of anthropogenic emissions affecting local-scale variation in water vapor. 
They also highlight the potential of stable isotopes of atmospheric vapor 
as a complementary tool for studying air pollution and its sources.












    